% Generated draft: requires mathematical and source review.
\documentclass[11pt]{article}
\usepackage[margin=25mm]{geometry}
\usepackage{fontspec}
\setmainfont{Latin Modern Roman}
\usepackage{amsmath,amssymb,mathtools,unicode-math}
\setmathfont{Latin Modern Math}
\usepackage{enumitem,needspace,xurl,hyperref}
\setlength{\parindent}{0pt}
\setlength{\parskip}{5pt}
\setlength{\emergencystretch}{3em}
\newcommand{\N}{\mathbb{N}}
\newcommand{\Z}{\mathbb{Z}}
\newcommand{\Q}{\mathbb{Q}}
\newcommand{\R}{\mathbb{R}}
\newcommand{\C}{\mathbb{C}}
\newcommand{\F}{\mathbb{F}}
\newcommand{\A}{\mathbb{A}}
\newcommand{\PP}{\mathbb{P}}
\newcommand{\E}{\mathbb{E}}
\newcommand{\eps}{\varepsilon}
\newcommand{\dd}{\,\mathrm{d}}
\begin{document}
\section*{UnsolvedMath ID: 20001276}
\textbf{Source problem code:} catalogue entry
\subsection{Definitions and mathematical statement}
The record merges two Section 9 problems and the next section's introductory heading. Let $K\subset\R^n$ be a convex body. For $a\in\operatorname{int}K$ define the (possibly asymmetric) Minkowski gauge $\gamma_{K-a}(x)=\inf\{t>0:x\in t(K-a)\}$ and the fixed-center vertex index
\[
V(K,a)=\inf_{\substack{m<\infty,\ p_1,\ldots,p_m\in\R^n\\K-a\subseteq\operatorname{conv}\{p_1,\ldots,p_m\}}}
\sum_{i=1}^m\gamma_{K-a}(p_i).
\]
For an origin-symmetric body the usual gauge is its norm and the centered invariant is $V(K,0)$. The first problem asks for the best universal upper bound in dimension three, equivalently the extremal value $M_3=\sup_{K=-K\subset\R^3}V(K,0)$. The attached aid records
\[
6\sqrt3\le M_3\le18,
\]
and a conjectured upper bound $12$, leaving the optimum open in its assessment.

The second source problem asks for an upper bound in general dimension with a center variable. The standard interior-centered version is $\operatorname{vein}_{\mathrm{int}}(K)=\inf_{a\in\operatorname{int}K}V(K,a)$. The aid states affine invariance, attainment of $V(K,a)$ for every fixed interior center after a finite cardinality reduction, the universal lower bound $V(K,a)\ge n+1$, and equality $V(K,a)=n+1$ for every simplex and every interior center. It does not supply the cardinality bound or a complete best dimension-dependent upper estimate. Fixed-center attainment does not assert attainment of the additional center infimum.

The printed formula instead permits $a\in K$, including boundary points, where the gauge need not be finite in every direction. Under an extended-gauge interpretation set $\gamma_{K-a}(0)=0$ and allow $+\infty$ elsewhere when appropriate. For a simplex $K=\operatorname{conv}\{v_0,\ldots,v_n\}$ at boundary center $a=v_0$, the containing configuration $0,v_1-v_0,\ldots,v_n-v_0$ has total cost $n$. Hence the boundary-allowed infimum is at most $n$ on a simplex, strictly below the interior-centered lower bound; no exact boundary infimum beyond this construction is asserted. The two conventions define materially different invariants and must not be silently interchanged.

The appended Section 10, ``Reconstruction of Polytopes with Few Facets (or Few Vertices),'' merely states motivation from efficient estimation using random points. It supplies no further reconstruction question in this record.
\subsection{Short English statement}
Find vertex-index upper bounds for symmetric three-dimensional bodies and general convex bodies. The supplied aid gives structural results but shows that the printed permission for boundary centers changes the invariant, even on simplices. The appended reconstruction heading is separate context.
\subsection{Sources}
\begin{itemize}
\item Local UnsolvedMath registry, permanent ID $20001276$, source code catalogue entry; canonical \texttt{aim-convex-geometry-notes.json}, \texttt{notes[7]}.
\item American Institute of Mathematics, 2013 workshop ``Sections of convex bodies,'' Section 9 problems 2 and 3, with a spillover into Section 10: \url{https://aimath.org/pastworkshops/sectionsconvexproblems.pdf}. The supplied verification identifies this partition of the official PDF; no PDF inspection is claimed here. No original references or remarks are supplied.
\item Attached research attempt 1: partial result, run and registry status partially solved, rigor proved, novelty confidence low, classification PARTIAL-PROGRESS. Its bounds and fixed-center claims are machine generated and not peer reviewed; independently verify before citation. No authoritative citations for the numeric bounds or attainment theorem are given.
\end{itemize}
\end{document}
